\begin{center}\begin{tikzpicture}[font=\small,>=stealth,line join=round]
\draw[gray!40,rounded corners=3pt] (0,0) rectangle (14.2,8.5);\node at (7.1,7.9) {在同一条竖线上，对称点必须重合};
\node at (3.5,6.85){假设 $c\ne0$：同一个输入出现两个值};\node at (10.6,6.85){$c=0$：两个点合为原点};\begin{scope}[shift={(3.5,3.8)},x=1cm,y=1cm]\draw[->,gray] (-2.1,0)--(2.1,0) node[right]{$x$};\draw[->,gray] (0,-2.0)--(0,2.0) node[above]{$y$};\fill[red] (0,1.3) circle(2.2pt) node[right=5pt] {$(0,c)$};\fill[red] (0,-1.3) circle(2.2pt) node[right=5pt] {$(0,-c)$};\draw[red,dashed] (0,-1.3)--(0,1.3);\end{scope}\begin{scope}[shift={(10.5,3.8)},x=1cm,y=1cm]\draw[->,gray] (-2.1,0)--(2.1,0) node[right]{$x$};\draw[->,gray] (0,-2)--(0,2) node[above]{$y$};\draw[blue,thick,domain=-1.2:1.2,samples=150] plot(\x,{(\x)^3});\fill[orange!85!black] (0,0) circle(2.2pt) node[below right=5pt] {$(0,0)$};\end{scope}\node at (3.5,1.7){左侧不能作为函数的图像};\node at (10.6,1.7){一个横坐标只对应一个纵坐标};
\node[text width=12.8cm,align=center] at (7.1,.65) {左图为 $c\ne0$ 的假设示意，并非函数图像；$c<0$ 时上下两点交换。};
\end{tikzpicture}\end{center}
